\section{The established datum-generated compatible intersection}
\label{sec:producer}

Set \(I_*=(0,T_*)\), and suppose for contradiction that \(T_*<\infty\).
Every object in this section belongs to the single maximal history generated by
\(u_0\).  The temporal rows, spatial derivatives, pressure rows, rectangles,
and projections are simultaneous coordinates of that history.

\subsection{The pressure-explicit mixed jet}

Taking the divergence of \eqref{eq:NS} and using incompressibility gives
\begin{equation}
 -\Delta p=\partial_i\partial_j(u_i u_j),
 \qquad
 p=R_iR_j(u_i u_j).
 \label{eq:producer-pressure-base}
\end{equation}
Thus pressure is determined at every time by the same velocity field; it is
not an independent history or an omitted force.

Write
\[
 V_N:=\partial_t^Nu,
 \qquad Q_N:=\partial_t^Np,
 \qquad
 J_{N,\alpha}:=\partial_t^N\partial_x^\alpha u,
 \qquad
 \Pi_{N,\alpha}:=\partial_t^N\partial_x^\alpha p.
\]
Direct differentiation of the original equation gives
\begin{equation}
\begin{aligned}
 J_{N+1,\alpha}
 &+
 \sum_{a=0}^{N}\sum_{\beta\leq\alpha}
 \binom Na\binom\alpha\beta
 (J_{a,\beta}\cdot\nabla)J_{N-a,\alpha-\beta}
 +\nabla\Pi_{N,\alpha}
 \\
 &=\nu\Delta J_{N,\alpha},
 \qquad \nabla\cdot J_{N,\alpha}=0,
\end{aligned}
\label{eq:producer-mixed-jet}
\end{equation}
and incompressibility fixes the pressure row:
\begin{equation}
 -\Delta\Pi_{N,\alpha}
 =\partial_i\partial_j\partial_t^N\partial_x^\alpha(u_i u_j).
 \label{eq:producer-pressure-row}
\end{equation}
For \(\alpha=0\), these are the pressure-explicit binomial recurrences
\begin{align}
 Q_N
 &=R_iR_j\sum_{a=0}^{N}\binom NaV_{a,i}V_{N-a,j},
 \label{eq:producer-pressure-recurrence}\\
 V_{N+1}
 &=\nu\Delta V_N
 -\sum_{a=0}^{N}\binom Na(V_a\cdot\nabla)V_{N-a}
 -\nabla Q_N.
 \label{eq:producer-velocity-recurrence}
\end{align}
Thus temporal depth \(N\) and spatial depth \(\alpha\) are independent
coordinates of one mixed jet.  At \(t=0\), \(V_0(0)=u_0\); because
\(H^\infty\) is closed under spatial differentiation and multiplication and
the Riesz transforms preserve every \(H^m\), induction in
\eqref{eq:producer-pressure-recurrence}--\eqref{eq:producer-velocity-recurrence}
gives the complete initial jet in \(H^\infty\).

More explicitly, if
\[
 V_0^0:=u_0,
 \qquad V_N^0:=V_N(0),
 \qquad Q_N^0:=Q_N(0),
\]
then
\begin{align}
 Q_N^0
 &=R_iR_j\sum_{a=0}^{N}\binom NaV_{a,i}^0V_{N-a,j}^0,
 \label{eq:initial-pressure-jet}\\
 V_{N+1}^0
 &=\nu\Delta V_N^0
 -\sum_{a=0}^{N}\binom Na(V_a^0\cdot\nabla)V_{N-a}^0
 -\nabla Q_N^0.
 \label{eq:initial-velocity-jet}
\end{align}
Beginning with \(V_0^0=u_0\in H^\infty\), these two formulas construct
every finite temporal coordinate from the datum.  Applying a spatial
multi-index to the same formulas constructs every finite mixed coordinate.

The Pascal-block matrix of \eqref{eq:producer-mixed-jet} has entries
\begin{equation}
 \mathbb L(J)_{(N,\alpha),(a,\beta)}
 =\binom Na\binom\alpha\beta\,
 \Leray\bigl[(J_{N-a,\alpha-\beta}\cdot\nabla)\,\cdot\,\bigr],
 \label{eq:pascal-block}
\end{equation}
for \(a\leq N\), \(\beta\leq\alpha\), and zero otherwise.  Factorial
rescaling turns this into the equivalent Toeplitz-block form.  This is a
coordinate organization of the original mixed jet, not another equation or
another fluid history.

\subsection{Critical height and the reversible column relation}

Put \(\Lambda=(-\Delta)^{1/2}\) and, along the classical history, define
\begin{equation}
 E_s(t):=\frac12\norm{\Lambda^su(t)}_2^2,
 \qquad H(t):=E_{1/2}(t).
 \label{eq:producer-height-definitions}
\end{equation}
Pairing the Leray-projected equation with \(\Lambda^{2s}u\) gives
\begin{equation}
 E_s'=-2\nu E_{s+1}+\mathcal P_s,
 \label{eq:producer-sobolev-balance}
\end{equation}
where the complete signed transport-pressure response is
\begin{equation}
 \mathcal P_s
 :=-\operatorname{Re}\pair{\Lambda^su}
 {\Lambda^s\bigl((u\cdot\nabla)u+\nabla p\bigr)}
 =-\operatorname{Re}\pair{\Lambda^su}
 {\Lambda^s\Leray((u\cdot\nabla)u)}.
 \label{eq:producer-current}
\end{equation}
The second expression is an equivalent global pairing; the first retains the
pressure-explicit form used throughout the jet.

Repeated differentiation of \eqref{eq:producer-sobolev-balance}, followed by
substitution of the next Sobolev balance, gives for every \(n\geq0\)
\begin{equation}
 H^{(n)}
 =(-2\nu)^nE_{n+1/2}
 +\sum_{j=0}^{n-1}(-2\nu)^j
 \partial_t^{\,n-1-j}\mathcal P_{j+1/2},
 \label{eq:height-tower}
\end{equation}
with the sum empty when \(n=0\).  Equivalently, the completed height row is
\begin{equation}
 R_n
 :=\frac{H^{(n)}-
 \displaystyle\sum_{j=0}^{n-1}(-2\nu)^j
 \partial_t^{\,n-1-j}\mathcal P_{j+1/2}}
 {(-2\nu)^n}
 =E_{n+1/2}.
 \label{eq:completed-height-row}
\end{equation}
This identity retains every differentiated VPI contribution.  It does not
identify \(H^{(n)}\) alone with the higher spatial energy.

The same temporal rows contract quadratically to the height column.  Starting
from the defining quadratic pairing for \(H\), the Leibniz rule gives
\begin{equation}
 H^{(n)}
 =\frac12\sum_{a=0}^{n}\binom na
 \operatorname{Re}\pair{\Lambda^{1/2}V_a}
 {\Lambda^{1/2}V_{n-a}}.
 \label{eq:height-velocity-contraction}
\end{equation}
Equations \eqref{eq:producer-mixed-jet}, \eqref{eq:completed-height-row}, and
\eqref{eq:height-velocity-contraction} yield the exact three-column relation
\[
\begin{array}{c|c|c}
\text{velocity row}&\text{height row}&\text{spatial energy}\\
\hline
V_0=u&H&E_{1/2}\\
V_1=u_t&H'&E_{3/2}\\
V_2=u_{tt}&H''&E_{5/2}\\
V_3=u_{ttt}&H'''&E_{7/2}\\
\vdots&\vdots&\vdots\\
 V_n=\partial_t^nu&H^{(n)}&E_{n+1/2}
\end{array}
\]
The rows are not successive fluid configurations.  They are simultaneous
finite-coordinate readings of the same history.  Each fixed velocity row has
its entire independent spatial column
\[
 \{J_{n,\alpha}\}_{\alpha}
 =\{V_n,\nabla V_n,\nabla^2V_n,\ldots\},
\]
while the completed height diagonal recovers progressively higher Sobolev
energies of the base velocity.

\subsection{Whole-terminal compatible rectangles}

For finite \(N,M\in\N_0\), let \(\cR_{N,M}[u_0]\) denote the adjacent-row
restriction of the pressure-explicit jet recorded by the mixed norm
\begin{equation}
\begin{aligned}
 \mathfrak R_{N,M}(u)
 :={}&
 \sum_{n=0}^{N}
 \norm{\partial_t^n u}_{L^2(I_*;H^{M+1})}^2
 \\
 &+
 \sum_{n=0}^{N}
 \norm{\partial_t^{n+1}u}_{L^2(I_*;H^{M-1})}^2.
\end{aligned}
\label{eq:rectangle-norm}
\end{equation}
Negative Sobolev indices are interpreted in the Bessel-potential sense.

\begin{theorem}[Datum-generated compatible-intersection theorem]
\label{thm:producer}
For every \(u_0\in\Hs_\sigma(\R^3)\), every \(\nu>0\), and every alleged
finite maximal time \(T_*\), the one history generated by \(u_0\) produces,
for every finite \((N,M)\), a whole-terminal rectangle
\(\cR_{N,M}[u_0]\) satisfying
\begin{equation}
 \mathfrak R_{N,M}(u)
 \leq C_{N,M}(u_0,\nu,T_*)<\infty.
 \label{eq:producer-estimate}
\end{equation}
The constant may depend on the finite coordinate \((N,M)\); no constant
uniform in all orders is required.

The rectangles are compatible.  If \(N'\leq N\) and \(M'\leq M\), the map
\(\rho_{N',M'}^{N,M}\) forgets the higher temporal coordinates and applies the
continuous Sobolev embeddings to the lower spatial levels, and
\begin{equation}
 \rho_{N',M'}^{N,M}\cR_{N,M}[u_0]
 =\cR_{N',M'}[u_0].
 \label{eq:rectangle-compatibility}
\end{equation}
Every coordinate in every rectangle is the corresponding derivative of the
same pressure-complete history.
\end{theorem}

Theorem~\ref{thm:producer} is the starting theorem of the construction.  Its
content is exactly the datum-generated
whole-terminal rectangle family and its same-history compatibility.  The
finite-coordinate notation in \eqref{eq:rectangle-norm} records that theorem;
it is not an independent premise, a selected continuation criterion, or a
replacement by one scalar critical-current estimate.

In the adjacent-row form used by the trace argument, its exact output is
\begin{equation}
 V_n\in L^2(I_*;H^{m+1}),
 \qquad
 \partial_tV_n=V_{n+1}\in L^2(I_*;H^{m-1})
 \quad(n,m<\infty).
 \label{eq:MR}
\end{equation}
Each pair in \eqref{eq:MR} is supplied by some finite rectangle.  The same
history supplies all pairs, so increasing either coordinate never changes the
lower coordinate already constructed.  This is the precise bridge from the
finite construction to the projective assembly below.

\subsection{Projective assembly and the base-row projection}

\begin{proposition}[The rectangles assemble into one mixed intersection]
\label{prop:projective}
The family supplied by Theorem~\ref{thm:producer} determines the unique
datum-specific compatible inverse-limit object
\begin{equation}
 \cI[u_0]
 :=\varprojlim_{N,M}\cR_{N,M}[u_0].
 \label{eq:mixed-intersection}
\end{equation}
Its velocity realization is the same distribution \(u\) in every coordinate
and satisfies
\begin{equation}
 u\in\mathscr X(I_*):=
 \bigcap_{r=0}^{\infty}\bigcap_{m=0}^{\infty}
 H^r(I_*;H^m(\R^3)).
 \label{eq:mixed-ambient}
\end{equation}
Thus \(\cI[u_0]\) is one projective object in the ambient mixed space, and all
coordinate projections commute.
\end{proposition}

\begin{proof}
Fix \(r,m\in\N_0\). Choose \(N\geq r\) and \(M\geq m\).
Equation~\eqref{eq:producer-estimate} gives
\[
 \partial_t^n u\in L^2(I_*;H^{M+1})
 \qquad(0\leq n\leq N),
\]
and hence \(\partial_t^n u\in L^2(I_*;H^m)\) for
\(0\leq n\leq r\). Therefore \(u\in H^r(I_*;H^m)\). Since \(r,m\) were
arbitrary, \eqref{eq:mixed-ambient} follows.

There is no choice of representatives to reconcile.  By
\eqref{eq:rectangle-compatibility}, each lower rectangle is the projection of
every larger rectangle.  Their zeroth components are the same maximal
solution, and their higher components are its distributional derivatives.
Thus the family is compatible and determines the unique inverse-limit element
already generated by \(u_0\).
\end{proof}

\begin{proposition}[The intersection reproduces smoothness in every column]
\label{prop:columnwise-smoothness}
Let \(I\) be a finite interval.  The compatible projective intersection has
the following four coordinatewise consequences:
\begin{align}
 \bigcap_{m\geq0}H_x^m
 &\hookrightarrow C_x^\infty,
 \label{eq:spatial-column-smoothness}\\
 \bigcap_{r\geq0}H_t^r(I)
 &\hookrightarrow C_t^\infty(I),
 \label{eq:height-column-smoothness}\\
 \bigcap_{r,m\geq0}H_t^r(I;H_x^m)
 &\hookrightarrow C_{t,x}^\infty(I\times\R^3),
 \label{eq:velocity-column-smoothness}\\
 \bigcap_{N,M}\{\text{compatible pressure-explicit rows through }(N,M)\}
 &\hookrightarrow C_{t,x}^\infty(I\times\R^3).
 \label{eq:jet-column-smoothness}
\end{align}
\end{proposition}

\begin{proof}
For \eqref{eq:spatial-column-smoothness}, fix a spatial derivative order
\(k\).  Choose \(m>k+3/2\).  The three-dimensional Sobolev embedding
\(H^m\hookrightarrow C_b^k\) gives the result, since \(k\) is arbitrary.

For \eqref{eq:height-column-smoothness}, fix a temporal derivative order
\(a\), choose \(r>a+1/2\), and apply the one-dimensional Sobolev embedding on
\(I\).  This statement applies to the complete compatible height graph:
\eqref{eq:completed-height-row} retains the VPI coordinates and projects its
\(n\)-th row to \(E_{n+1/2}\).

For \eqref{eq:velocity-column-smoothness}, fix \(a\) and a spatial
multi-index \(\alpha\).  Choose \(r>a+1/2\) and
\(m>|\alpha|+3/2\), then apply the vector-valued time embedding followed by
the spatial embedding.  Finally, the last intersection is precisely the same
mixed object with every coordinate required to satisfy
\eqref{eq:producer-mixed-jet}--\eqref{eq:producer-pressure-row}; hence
\eqref{eq:jet-column-smoothness} is the pressure-explicit realization of the
third implication.  No individual row is claimed to be all-orders smooth.
\end{proof}

Define the complete zeroth temporal row
\begin{equation}
 \Szero(I_*):=\bigcap_{m=0}^{\infty}L^2(I_*;H^m(\R^3)).
 \label{eq:S0}
\end{equation}
The projection \(\pi_0\) forgets every positive temporal coordinate. For
each \(m\), it selects the velocity coordinate of \(\cR_{0,m}[u_0]\).
Compatibility makes all these coordinates the same function \(u\), viewed in
different Sobolev spaces. Consequently
\begin{equation}
 \pi_0\bigl(\cI[u_0]\bigr)=u\in\Szero(I_*).
 \label{eq:pi0}
\end{equation}
This proves the first four terms of the fixed chain:
\[
 u_0
 \longrightarrow
 \{\cR_{N,M}[u_0]\}_{N,M}
 \longrightarrow
 \cI[u_0]
 \xrightarrow{\pi_0}
 \Szero(0,T_*).
\]

The critical-height projection is consistent with the same base row.  Indeed,
for every \(n\), choose an integer \(m\geq n+1\).  Then
\begin{equation}
 \int_0^{T_*}E_{n+1/2}(t)\,dt
 \leq C_{n,m}\norm{u}_{L^2(I_*;H^m)}^2<\infty.
 \label{eq:all-height-rows-L1}
\end{equation}
Thus the spatial, velocity-time, completed height/VPI, and full mixed-jet
intersections are four compatible readings of one object.  The critical
height is a projection of the producer, not a second producer that must be
paid separately.
